\chapter{Robustesse : la dégradation honnête}
\label{ch:robustesse}

Un service qui dépend de six sources publiques tombe souvent — non pas
complètement, mais \emph{partiellement}. La qualité de l'architecture se mesure
donc moins à son fonctionnement nominal qu'à sa façon de se dégrader.

\section{Diagramme d'activité : collecter sans mentir}

\begin{figure}[h]
\centering
\figbox{
\begin{tikzpicture}[font=\scriptsize,
    act/.style={draw=typhoonBlue!75,fill=typhoonBlue!5,rounded corners=2pt,
                align=center,minimum height=0.75cm,inner sep=3pt},
    bar/.style={fill=black!75,minimum height=0.09cm},
    dec/.style={diamond,draw=typhoonAmber!90,fill=typhoonAmber!8,aspect=2.2,
                align=center,inner sep=1pt}]

  \node[circle,fill=black!75,minimum size=3mm,inner sep=0pt] (start) at (4.5,0.4) {};
  \node[bar,minimum width=7.8cm] (fork) at (4.5,-0.35) {};

  \node[act,minimum width=1.85cm] (a1) at (0.9,-1.3) {BDNB};
  \node[act,minimum width=1.85cm] (a2) at (2.9,-1.3) {Géorisques};
  \node[act,minimum width=1.85cm] (a3) at (4.9,-1.3) {TRI au point};
  \node[act,minimum width=1.85cm] (a4) at (6.9,-1.3) {Hydro, météo};
  \node[act,minimum width=2.2cm] (a5) at (3.9,-2.25) {Photo de terrain};

  \node[bar,minimum width=7.8cm] (join) at (4.5,-2.95) {};

  \node[dec,minimum width=2.5cm,minimum height=1.0cm] (d) at (4.5,-3.85) {source\\muette ?};

  \node[act,minimum width=3.5cm,fill=typhoonRed!6,draw=typhoonRed!70] (abs) at (1.7,-5.15)
    {absence déclarée\\\scriptsize avec sa raison nommée};
  \node[act,minimum width=3.5cm,fill=typhoonGreen!7,draw=typhoonGreen!70] (fait) at (7.3,-5.15)
    {fait publié\\\scriptsize source + résolution + attribution};

  \node[act,minimum width=6.4cm] (asm) at (4.5,-6.35) {assembler et valider le contrat};
  \node[circle,draw=black!75,thick,minimum size=4.2mm,inner sep=0pt] (end) at (4.5,-7.15) {};
  \node[circle,fill=black!75,minimum size=2.4mm,inner sep=0pt] at (4.5,-7.15) {};

  \draw[flow] (start) -- (fork);
  \foreach \a in {a1,a2,a3,a4} {\draw[flow] (fork.south -| \a.north) -- (\a.north);}
  \draw[flow] (a2.south) -- (a5.north);
  \foreach \a in {a1,a2,a3,a4,a5} {\draw[flow] (\a.south) -- (\a.south |- join.north);}
  \draw[flow] (join) -- (d);
  \draw[flow] (d.west) -- node[above]{oui} (abs.north);
  \draw[flow] (d.east) -- node[above]{non} (fait.north);
  \draw[flow] (abs.south) -- (abs.south |- asm.north);
  \draw[flow] (fait.south) -- (fait.south |- asm.north);
  \draw[flow] (asm) -- (end);
\end{tikzpicture}
}
\caption{Diagramme d'activité de la collecte. La branche « source muette »
produit un \emph{fait} sur la source, pas une valeur sur le risque : c'est cette
distinction qui rend la dégradation exploitable.}
\label{fig:activite}
\end{figure}

\section{Diagramme d'états : les cinq vérités de l'absence}
\label{sec:etats-tri}

Le cas du périmètre TRI est le plus instructif, parce que « il n'y a pas de
classe de crue ici » recouvre en réalité cinq situations différentes. Les
confondre revient à présenter une panne comme une bonne nouvelle.

\begin{figure}[h]
\centering
\figbox{
\begin{tikzpicture}[font=\scriptsize,
    st/.style={draw=typhoonBlue!75,fill=typhoonBlue!5,rounded corners=9pt,
               align=center,minimum width=3.1cm,minimum height=0.8cm,inner sep=3pt},
    stok/.style={st,draw=typhoonGreen!75,fill=typhoonGreen!7},
    stwarn/.style={st,draw=typhoonRed!75,fill=typhoonRed!6}]

  \node[circle,fill=black!75,minimum size=3mm,inner sep=0pt] (init) at (0,1.5) {};
  \node[st] (q) at (0,0.6) {Interrogation du\\périmètre TRI};

  \node[stok] (hors) at (-2.6,-1.3) {\textbf{Hors TRI}\\fait territorial};
  \node[st] (sans) at (2.6,-1.3) {\textbf{Dans un TRI}\\sans classe au point};
  \node[stwarn] (nu) at (-2.6,-3.1) {\textbf{Appartenance}\\non vérifiée};
  \node[stwarn] (panne) at (2.6,-3.1) {\textbf{Service}\\indisponible};

  \draw[flow] (init) -- (q);
  \draw[flow] (q.south) -- ++(0,-0.25) -| (hors.north);
  \draw[flow] (q.south) -- ++(0,-0.25) -| (sans.north);
  \draw[flow] (q.south) -- ++(0,-0.5) -| (nu.north);
  \draw[flow] (q.south) -- ++(0,-0.5) -| (panne.north);

  \node[font=\tiny,anchor=east,align=right] at (-0.15,-0.62) {point hors emprise};
  \node[font=\tiny,anchor=west,align=left] at (0.15,-0.62) {polygone couvrant,\\ aucune hauteur publiée};
  \node[font=\tiny,anchor=east,align=right] at (-0.4,-1.95) {réponse inexploitable};
  \node[font=\tiny,anchor=west,align=left] at (0.4,-1.95) {erreur réseau, 5xx, quota};

  \node[umlnote,below=0.35cm of nu,text width=8.6cm,anchor=north]
    {Chaque état a sa \emph{formulation propre}, partagée entre le panneau de
     scénarios et le rapport (« ce n'est pas une absence de risque »). Aucun
     chemin ne produit « 0 m » à la place d'une source muette.};
\end{tikzpicture}
}
\caption{Diagramme d'états de l'absence TRI. « Hors TRI » est un fait ; « service
indisponible » est une panne : deux messages, deux traitements.}
\label{fig:etats-tri}
\end{figure}

\begin{invariant}[title={Absence \(\neq\) panne}]
Le connecteur porte cinq vérités distinctes : chargement, panne, dans un TRI sans
classe au point, hors TRI, appartenance non vérifiée. La formulation est unique et
partagée. Un panneau qui afficherait « pas de risque » sur une panne violerait la
règle la plus importante du projet.
\end{invariant}

\section{La discipline de test}

La robustesse n'est pas vérifiée par relecture mais par exécution. La suite
offline couvre 261 tests et ne dépend d'aucun accès réseau : toutes les sources
sont simulées par un transport HTTP factice, ce qui rend les scénarios de panne
\textbf{reproductibles} — un 500, un 429, un GML tronqué se rejouent à volonté.

\begin{itemize}
  \item \textbf{Les tests précèdent le code.} Un test qui échoue est commité
    avant l'implémentation qui le fait passer.
  \item \textbf{Les tests existants ne sont pas modifiés} pour faire passer une
    fonctionnalité : ce serait déplacer la barre au lieu de franchir l'obstacle.
  \item \textbf{Un garde-fou d'isolation} vérifie qu'aucun module factuel
    n'importe le bac à sable, et que le bac à sable n'importe ni le moteur de
    dommages, ni le rapport, ni la carte, ni le schéma du contrat.
  \item \textbf{La limitation de débit est testée} sous, à et au-delà du seuil,
    y compris le fait que les routes de lecture légère ne sont jamais freinées.
\end{itemize}

\section{Ce que la dégradation produit, en pratique}

Lors de la vérification en conditions réelles, deux sources ont effectivement
échoué pendant la même exécution : la BDNB a répondu un statut de limitation sur
sa première page, et le service TRI s'est déclaré partiellement indisponible.

Le résultat obtenu est exactement celui attendu par la figure~\ref{fig:activite} :

\begin{itemize}
  \item la fiche bâtiment a affiché des champs vides plutôt que des valeurs
    inventées ;
  \item le panneau TRI a écrit « point non cartographié » et rappelé que ce
    n'était pas une absence de risque ;
  \item le rapport a signalé douze aléas « non vérifiés » ;
  \item \textbf{aucune erreur serveur n'a été produite} : le diagnostic a abouti,
    partiellement renseigné et entièrement honnête.
\end{itemize}

\begin{encadre}[title={Le critère de réussite n'est pas le taux de succès}]
Un diagnostic dégradé mais exact vaut mieux qu'un diagnostic complet et faux. Le
service préfère rendre 80 \% de faits vérifiés avec leurs manques déclarés,
plutôt que 100 \% de champs remplis dont une partie serait devinée.
\end{encadre}
